%!TEX root = ../../../main.tex

\chapter{Reflexion und Ausblick}
\label{chap:06-reflection}

% Bewertungskriterium: Literaturbezug und Reflexion (10 %), Budget 0,75 Seiten
% Belege: evaluation-log (RUN-010 bis RUN-021), ADR-0013/0015/0016, ISS-062/067/070/071/074

Die Annahme aus Kapitel~\ref{chap:01-problem-use-case}, dass Prüffragen einen Agenten brauchen, hat sich in diesem Prototyp nicht bestätigt. Auch bei den vier Grenzfällen war P2 nicht besser als die Pipeline P1 (keine richtige Antwort gegen eine), der Agent brauchte je Frage aber fast dreimal so lange. Ferrazzi et al. fanden ebenfalls keinen generellen Vorteil von agentischem RAG gegenüber einer ausgebauten Pipeline, wohl aber im Mittel 3,3-mal so viele Eingabe-Token~\cite{ferrazzi2026agentic}. Gegen agentisches RAG an sich spricht das nicht. P2 hatte mehr Sollfakten im Kontext als P1, scheiterte aber am Lesen und Vergleichen der Belege, also vor allem an den Grenzen eines 7B-Modells auf der CPU. Mit einem stärkeren Modell oder mit Fragen, deren Suchweg nicht vorab feststeht, kann der Vergleich anders ausfallen.

Rückblickend hätte jede Änderung am Agenten einen eigenen Lauf gebraucht. Nach dem Umbau der Suchwerkzeuge stieg P2 mit \texttt{llama3.1:8b} von 3 auf 9 richtige Antworten. Danach kamen Wiederholungssperre, Rechennachfrage und neue Suchregeln im Prompt auf einmal dazu, und P2 fiel auf 6. Eine Wiederholung ohne Codeänderung war wortgleich, der Rückgang lag also am Code. Welche der drei Änderungen ihn auslöste, ist offen, denn bei 60 bis 90 Minuten je Lauf wurden sie zusammen gemessen.

Mehr Prüfungen im Code machten die Antworten nicht richtiger. Einen eingebauten Fehler fand der Agent mit tragfähiger Begründung nur in einem Zwischenstand (G4, 95,78\,EUR Abweichung). Ein gültiger Aufruf von \texttt{calculate} belegt keine richtige Rechnung, bei G2 zog das Modell einen falschen Istwert vom Sollwert ab und meldete einen Fehlalarm. Die Rückfragen prüfen nur, ob eine Belegnummer oder ein Rechenaufruf vorkommt, nicht ob der Beleg die Aussage stützt. Auch der feste Prüfablauf sicherte nur die Herkunft jeder Zahl. In 4 von 14 Prüfgrößen stammten Soll- und Istwert aus demselben Treffer, und nach seiner zusätzlichen Nachricht antwortete das Modell bei S5 auf Chinesisch.

Ein nächster Versuch sucht Soll- und Istwert getrennt nach Dokumentart und prüft, ob beide aus verschiedenen Treffern stammen und dieselbe Größe meinen. Ob das wirkt, zeigt erst ein größerer Fragenkatalog mit einzelnen, wiederholten Messungen.
